\documentclass[a4paper,12pt]{article}
\usepackage[ngerman]{babel}
\usepackage[utf8]{inputenc}
\usepackage[T1]{fontenc}
\usepackage{amsmath,amssymb}
\usepackage{tikz}
\usetikzlibrary{automata,positioning,arrows.meta}

\tikzset{
    ->,>=Stealth,
    node distance=2.5cm,
    every state/.style={thick,minimum size=1cm},
    initial text={},
}

\newcommand{\defbox}[1]{%
\medskip
\noindent\fbox{%
    \begin{minipage}{0.94\linewidth}
    \textbf{Definition / Vorgehen.}\par\smallskip
    #1
    \end{minipage}%
}
\medskip\par
}

\title{ETI -- Aufgabenblatt 2 \\ Hausaufgaben}
\author{}
\date{}

\begin{document}
\raggedright
\setlength{\parindent}{0pt}

{\LARGE\bfseries ETI -- Aufgabenblatt 2 \\ Hausaufgaben\par}
\vspace{2em}

\section*{Aufgabe 2.5}

\defbox{
Ein NFA akzeptiert ein Wort, wenn mindestens ein möglicher Berechnungspfad in einem akzeptierenden Zustand endet.
Die Sprache
\[
B=\{a,b\}^* \circ \{a\} \circ \{a,b\}^2
\]
besteht genau aus den Wörtern über $\{a,b\}$, deren drittes Zeichen von rechts ein $a$ ist. Für den NFA kann man dieses $a$ nichtdeterministisch raten. Für den DFA muss man deterministisch die letzten drei gelesenen Zeichen speichern.
}

\subsection*{1. NFA $M_1$ mit 4 Zuständen}

Der NFA rät beim Lesen eines $a$, dass dieses $a$ das drittletzte Zeichen ist, und liest danach noch genau zwei Zeichen.

\[
M_1=(Q,\Sigma,\delta,q_0,F)
\]
mit
\[
Q=\{q_0,q_1,q_2,q_3\}, \quad \Sigma=\{a,b\}, \quad F=\{q_3\}.
\]
Die Übergänge sind:
\[
\delta(q_0,a)=\{q_0,q_1\}, \quad
\delta(q_0,b)=\{q_0\},
\]
\[
\delta(q_1,a)=\delta(q_1,b)=\{q_2\}, \quad
\delta(q_2,a)=\delta(q_2,b)=\{q_3\},
\]
alle anderen Übergänge führen in die leere Menge.

\begin{flushleft}
\begin{tikzpicture}[node distance=2.8cm]
    \node[state, initial] (q0) {$q_0$};
    \node[state, right=of q0] (q1) {$q_1$};
    \node[state, right=of q1] (q2) {$q_2$};
    \node[state, accepting, right=of q2] (q3) {$q_3$};

    \path (q0) edge[loop above] node {$a,b$} ()
          (q0) edge node[above] {$a$} (q1)
          (q1) edge node[above] {$a,b$} (q2)
          (q2) edge node[above] {$a,b$} (q3);
\end{tikzpicture}
\end{flushleft}

\textbf{Korrektheit:}
Der Automat kann beliebig lange in $q_0$ bleiben. Sobald er ein $a$ liest, kann er nach $q_1$ wechseln und dieses $a$ als drittletztes Zeichen raten. Danach müssen noch genau zwei Zeichen gelesen werden, um in $q_3$ zu enden. Also akzeptiert $M_1$ genau die Wörter aus $\{a,b\}^*a\{a,b\}^2=B$.

\subsection*{2. DFA $M_2$ mit 8 Zuständen}

Der DFA speichert immer die letzten drei Zeichen. Am Anfang werden fehlende Zeichen links mit $b$ aufgefüllt, damit Wörter der Länge $0,1,2$ nicht versehentlich akzeptiert werden.

\[
M_2=(Q,\Sigma,\delta,bbb,F)
\]
mit
\[
Q=\{a,b\}^3, \quad \Sigma=\{a,b\},
\]
\[
\delta(xyz,\sigma)=yz\sigma
\quad\text{für }x,y,z,\sigma\in\{a,b\},
\]
und
\[
F=\{aaa,aab,aba,abb\}.
\]

\textbf{Übergangstabelle:}
\[
\begin{array}{c|cc}
    & a & b\\ \hline
aaa & aaa & aab\\
aab & aba & abb\\
aba & baa & bab\\
abb & bba & bbb\\
baa & aaa & aab\\
bab & aba & abb\\
bba & baa & bab\\
bbb & bba & bbb
\end{array}
\]

\textbf{Korrektheit:}
Nach jedem gelesenen Wort $w$ enthält der Zustand die letzten drei Zeichen von $w$, wobei am Anfang links mit $b$ aufgefüllt wird. Ein Zustand ist genau dann akzeptierend, wenn sein erstes Zeichen $a$ ist. Das ist genau die Aussage, dass das drittletzte Zeichen des Eingabewortes ein $a$ ist. Also gilt $L(M_2)=B$.

\section*{Aufgabe 2.6}

\defbox{
Das Schubfachprinzip besagt: Werden $8$ Objekte auf höchstens $7$ Fächer verteilt, dann liegen zwei verschiedene Objekte im selben Fach. Für einen DFA bedeutet das: Wenn $8$ verschiedene Wörter nach höchstens $7$ Zuständen führen, dann müssen zwei verschiedene Wörter denselben Zustand erreichen. Wörter, die denselben Zustand erreichen, können durch dieselbe Fortsetzung nicht mehr unterschieden werden.
}

Wir beweisen per Widerspruch. Angenommen, es gibt einen DFA
\[
M=(Q,\{a,b\},\delta,q_0,F)
\]
mit $|Q|=7$ und $L(M)=B$.

Betrachte die $8$ Wörter aus $\{a,b\}^3$. Da $M$ nur $7$ Zustände hat, gibt es nach dem Schubfachprinzip zwei verschiedene Wörter
\[
x=x_1x_2x_3,\qquad y=y_1y_2y_3
\]
aus $\{a,b\}^3$, die nach dem Einlesen im selben Zustand landen:
\[
\delta^*(q_0,x)=\delta^*(q_0,y).
\]

Da $x\neq y$, unterscheiden sich die Wörter an mindestens einer Position.

\begin{itemize}
    \item Falls $x_1\neq y_1$, dann liegt genau eines von $x,y$ in $B$, denn bei Wörtern der Länge $3$ ist das erste Zeichen das drittletzte Zeichen.
    \item Falls $x_2\neq y_2$, dann liegt genau eines von $xa,ya$ in $B$, denn in $xa$ bzw. $ya$ ist das zweite Zeichen von $x$ bzw. $y$ das drittletzte Zeichen.
    \item Falls $x_3\neq y_3$, dann liegt genau eines von $xaa,yaa$ in $B$, denn in $xaa$ bzw. $yaa$ ist das dritte Zeichen von $x$ bzw. $y$ das drittletzte Zeichen.
\end{itemize}

In jedem Fall gibt es also eine Fortsetzung $z\in\{\varepsilon,a,aa\}$, so dass genau eines der Wörter $xz$ und $yz$ in $B$ liegt.
Weil $x$ und $y$ aber denselben Zustand erreichen, gilt
\[
\delta^*(q_0,xz)=\delta^*(q_0,yz).
\]
Der DFA müsste also beide Wörter gleich behandeln. Das widerspricht $L(M)=B$.

Damit kann es keinen DFA mit $7$ Zuständen geben, der $B$ erkennt.

\section*{Aufgabe 2.7}

\defbox{
Ein NFA akzeptiert ein Wort, wenn es mindestens einen Pfad vom Startzustand zu einem akzeptierenden Zustand gibt, dessen Kantenbeschriftungen genau das Wort ergeben. $\varepsilon$-Übergänge verbrauchen kein Eingabezeichen. Um die Sprache zu bestimmen, zerlegt man die möglichen akzeptierenden Pfade in reguläre Ausdrücke.
}

Sei
\[
\Sigma=\{a,b,c\}.
\]
Der Startzustand ist $q_0$, die akzeptierenden Zustände sind $a_3,b_3,c_1$.

Vom Zustand $q_0$ aus kann der Automat ohne endgültiges Akzeptieren wieder nach $q_0$ zurückkehren, indem er entweder
\[
a
\]
liest, oder indem er
\[
c(ab)^*
\]
liest. Setze daher
\[
R = a \cup c(ab)^*.
\]
Jeder akzeptierende Pfad besteht aus beliebig vielen solchen Rückkehr-Blöcken und endet dann in einem der akzeptierenden Teile.

\medskip
\textbf{Ende in $c_1$:}
Der Pfad
\[
q_0 \xrightarrow{a} c_1
\]
liest das Wort $a$.

\medskip
\textbf{Ende in $a_3$:}
Der obere Teil liest
\[
c(ab)^+.
\]
Denn von $q_0$ geht es mit $c$ nach $a_1$, anschließend muss mindestens einmal $ab$ gelesen werden, um in den akzeptierenden Zustand $a_3$ zu gelangen.

\medskip
\textbf{Ende in $b_3$:}
Über den mittleren Teil kann der Automat zunächst beliebig viele Zeichen in $b_1$ lesen. Dann liest er ein $b$, danach beliebig viele weitere $b$ und anschließend ein beliebiges Zeichen aus $\Sigma$, um $b_3$ zu erreichen. Falls danach weitergelesen werden soll, muss mit $a$ oder $c$ nach $b_1$ zurückgegangen werden. Der dadurch beschriebene Ausdruck ist
\[
S=\Sigma^*bb^*\Sigma\bigl((a\cup c)\Sigma^*bb^*\Sigma\bigr)^*.
\]

Damit ist die akzeptierte Sprache
\[
\boxed{
L
=
R^*\bigl(a \cup c(ab)^+ \cup S\bigr)
}
\]
mit
\[
R=a\cup c(ab)^*
\quad\text{und}\quad
S=\Sigma^*bb^*\Sigma\bigl((a\cup c)\Sigma^*bb^*\Sigma\bigr)^*.
\]

\textbf{Begründung:}
Die Faktoren aus $R^*$ beschreiben genau die Pfade, die wieder bei $q_0$ enden. Danach muss ein akzeptierender Zustand erreicht werden. Das geschieht entweder durch den einzelnen Übergang nach $c_1$, durch den oberen Teil nach $a_3$, oder durch den mittleren Teil nach $b_3$. Diese drei Möglichkeiten sind genau die drei Terme $a$, $c(ab)^+$ und $S$.

\section*{Aufgabe 2.8}

\defbox{
Eine Sprache heißt regulär, wenn sie von einem DFA oder äquivalent von einem NFA erkannt wird. Für das Umkehren einer regulären Sprache verwendet man einen NFA: Man dreht alle Kanten um, macht die alten Endzustände zu möglichen Startzuständen und den alten Startzustand zum einzigen Endzustand.
}

Zu zeigen ist:
\[
L \text{ ist regulär}
\quad\Longleftrightarrow\quad
L^R \text{ ist regulär}.
\]

\subsection*{``$\Rightarrow$''}

Sei $L$ regulär. Dann gibt es einen NFA
\[
N=(Q,\Sigma,\delta,q_0,F)
\]
mit $L(N)=L$.

Wir konstruieren einen NFA
\[
N^R=(Q\cup\{q_s\},\Sigma,\delta_R,q_s,\{q_0\})
\]
für $L^R$. Dabei ist $q_s$ ein neuer Startzustand.

Die Übergänge von $N^R$ sind:
\begin{itemize}
    \item Für jeden alten Endzustand $f\in F$ gibt es einen $\varepsilon$-Übergang
    \[
    q_s \xrightarrow{\varepsilon} f.
    \]
    \item Für jeden Übergang $p\xrightarrow{x}q$ in $N$ mit $x\in\Sigma\cup\{\varepsilon\}$ gibt es in $N^R$ den umgedrehten Übergang
    \[
    q\xrightarrow{x}p.
    \]
\end{itemize}

Ein akzeptierender Pfad in $N$ mit Beschriftung
\[
w=w_1w_2\cdots w_n
\]
von $q_0$ zu einem Zustand $f\in F$ wird dadurch zu einem Pfad in $N^R$ von $q_s$ über $f$ nach $q_0$ mit Beschriftung
\[
w_n\cdots w_2w_1=w^R.
\]
Also erkennt $N^R$ genau $L^R$. Damit ist $L^R$ regulär.

\subsection*{``$\Leftarrow$''}

Sei nun $L^R$ regulär. Nach der gerade bewiesenen Richtung ist dann auch
\[
(L^R)^R
\]
regulär. Für jedes Wort gilt $(w^R)^R=w$, also ist
\[
(L^R)^R=L.
\]
Damit ist auch $L$ regulär.

\medskip
Folglich gilt: $L$ ist genau dann regulär, wenn $L^R$ regulär ist.

\end{document}
